\documentclass[11pt,a4paper]{article}
\usepackage[utf8]{inputenc}
\usepackage[T1]{fontenc}
\usepackage[ngerman]{babel}
\usepackage[a4paper,left=3cm,right=2cm,top=3cm,bottom=2cm]{geometry}
\usepackage{tikz}
\usetikzlibrary{positioning,arrows.meta,calc}
\usepackage{lmodern}
\usepackage{microtype}
\usepackage{booktabs}
\usepackage{array}
\usepackage{tabularx}
\usepackage{ragged2e}
\setlength{\parindent}{0pt}
\setlength{\parskip}{0.5\baselineskip}
\pagestyle{empty}

% Prevent hyphenation in table cells
\newcolumntype{L}[1]{>{\RaggedRight\hspace{0pt}}p{#1}}

\begin{document}

% ── Titelblock ───────────────────────────────────────────────────────────────
\begin{center}
  {\Large\bfseries Gliederung und Zeitplan}\\[0.2em]
  {\small Stand: 22. April 2026}
\end{center}

\noindent\rule{\linewidth}{0.4pt}\vspace{0.3em}

\noindent\begin{tabular}{@{}p{3.5cm} p{12cm}@{}}
\textbf{Titel}         & Ein Mehrkriterien-Entscheidungsmodell zur Auswahl eines KPI-Hub-Integrationsansatzes in einer fragmentierten Enterprise-Systemlandschaft \\[0.3em]
\textbf{Verfasser}     & Haakon Bessert (Matrikelnr. 3005149) \\
\textbf{Praxispartner} & SAP SE --- Janine Steidelmüller \\
\textbf{Betreuer DHSN} & Prof.\ Dr.-Ing.\ Jürgen Sachse \\
\textbf{Abgabe}        & 10. Juli 2026 \\
\end{tabular}

\vspace{0.3em}\noindent\rule{\linewidth}{0.4pt}

% ── Kapitelübersicht ─────────────────────────────────────────────────────────
{\large\bfseries Gliederung}\vspace{0.4em}

\noindent Die Arbeit umfasst 10 Kapitel auf 62 Seiten. Der detaillierte Aufbau ist nachfolgend dargestellt.

\vspace{0.3em}\begin{tabular}{@{}r p{12.5cm} r@{}}
\toprule
\textbf{Nr.} & \textbf{Kapitel} & \textbf{S.} \\
\midrule
1  & Einleitung & 4 \\
2  & Grundlagen und Stand der Forschung & 9 \\
3  & Methodik & 6 \\
4  & Ausgangssituation und Systemlandschaft & 7 \\
5  & Anforderungsanalyse & 7 \\
6  & Lösungsraum und Architektur-/Tooloptionen & 6 \\
7  & Entscheidungs- und Bewertungsmodell & 8 \\
8  & MVP: Design, Umsetzung und Validierung & 8 \\
9  & Ergebnisse und Diskussion & 5 \\
10 & Fazit und Ausblick & 2 \\
\midrule
   & \textbf{Gesamt} & \textbf{62} \\
\bottomrule
\end{tabular}

% ── Detaillierte Gliederung ───────────────────────────────────────────────────
\vspace{0.4em}
{\normalsize\bfseries Detaillierte Untergliederung}\vspace{0.3em}

\renewcommand{\arraystretch}{1.15}
\begin{tabular}{@{}L{7.5cm} L{8cm}@{}}
  % Kap 1
  \textbf{1\enspace Einleitung}                          & \textbf{6\enspace Lösungsraum und Architektur-/Tooloptionen} \\
  {\small 1.1\enspace Ausgangssituation und Problemstellung} & {\small 6.1\enspace Identifikation relevanter Lösungsansätze} \\
  {\small 1.2\enspace Lösungsrelevanz und Motivation}       & {\small 6.2\enspace Portal-/Dashboard-Föderation (Embedding)} \\
  {\small 1.3\enspace Einordnung in die Wirtschaftsinformatik} & {\small 6.3\enspace Zentrale BI-Aggregation} \\
  {\small 1.4\enspace Zielsetzung und Forschungsfrage}       & {\small 6.4\enspace Low-Code / No-Code Plattformen} \\
  {\small 1.5\enspace Aufbau der Arbeit}                     & {\small 6.5\enspace Kundenspezifische Webanwendung} \\
  & {\small 6.6\enspace Referenzprojekte bei SAP} \\
  \textbf{2\enspace Grundlagen und Stand der Forschung}    & {\small 6.7\enspace Vergleichende Übersicht und Machbarkeitseinschätzung} \\
  {\small 2.1\enspace KPI-Management und Business Intelligence} & \\
  {\small 2.2\enspace Enterprise-Systemintegration}          & \textbf{7\enspace Entscheidungs- und Bewertungsmodell} \\
  {\small 2.3\enspace Data Governance und Zugriffskontrolle} & {\small 7.1\enspace Wahl der Bewertungsmethodik} \\
  {\small 2.4\enspace Entscheidungsmethoden im Software Engineering} & {\small 7.2\enspace Aufbau des Scoring-Modells} \\
  {\small 2.5\enspace Defizite bestehender Ansätze}          & {\small 7.3\enspace Bewertung der Lösungsoptionen} \\
                                                             & {\small 7.4\enspace Ergebnis und Empfehlung} \\
  \textbf{3\enspace Methodik}                              & {\small 7.5\enspace Sensitivitätsanalyse} \\
  {\small 3.1\enspace Forschungsansatz (Design Science Research)} & \\
  {\small 3.2\enspace Literaturrecherche}                    & \textbf{8\enspace MVP: Design, Umsetzung und Validierung} \\
  {\small 3.3\enspace Ergebniserwartung}                     & {\small 8.1\enspace MVP-Zielsetzung und Scope} \\
  {\small 3.4\enspace Vorgehen im Überblick}                 & {\small 8.2\enspace Architektur und Design} \\
                                                             & {\small 8.3\enspace Umsetzung} \\
  \textbf{4\enspace Ausgangssituation und Systemlandschaft} & {\small 8.4\enspace Validierung} \\
  {\small 4.1\enspace Organisatorischer Kontext}             & \\
  {\small 4.2\enspace Bestehende Systemlandschaft}           & \textbf{9\enspace Ergebnisse und Diskussion} \\
  {\small 4.3\enspace Ist-Zustand der KPI-Versorgung}        & {\small 9.1\enspace Ergebnisse der Bewertung} \\
  {\small 4.4\enspace Datenbereitstellung und Schnittstellenlandschaft} & {\small 9.2\enspace Beantwortung der Forschungsfrage} \\
  {\small 4.5\enspace Identifizierte Problemfelder}          & {\small 9.3\enspace Kritische Diskussion} \\
                                                             & {\small 9.4\enspace Vergleich mit Ergebniserwartung} \\
  \textbf{5\enspace Anforderungsanalyse}                   & \\
  {\small 5.1\enspace Vorgehen bei der Anforderungserhebung} & \textbf{10\enspace Fazit und Ausblick} \\
  {\small 5.2\enspace Stakeholder und Use Cases}             & {\small 10.1\enspace Fazit} \\
  {\small 5.3\enspace Funktionale Anforderungen}             & {\small 10.2\enspace Beitrag der Arbeit} \\
  {\small 5.4\enspace Nicht-funktionale Anforderungen}       & {\small 10.3\enspace Ausblick} \\
  {\small 5.5\enspace Priorisierung der Anforderungen}       & \\
\end{tabular}

\vspace{0.5em}\noindent\rule{\linewidth}{0.4pt}

\newpage
% ── Roter Faden ───────────────────────────────────────────────────────────────
{\large\bfseries Roter Faden}\vspace{0.4em}

\noindent Die Kapitelstruktur folgt dem DSR-Prozessmodell nach Peffers et al.\ (2007) --- von der Problemdefinition über die Artefaktkonstruktion bis zur Evaluation:

\vspace{0.6em}
\begin{tikzpicture}[
  box/.style={
    rectangle, draw, rounded corners=3pt,
    minimum height=1.6cm, minimum width=2.7cm, text width=2.5cm,
    align=center, font=\small
  },
  arr/.style={-{Stealth[scale=1.1]}, thick},
  phase/.style={font=\scriptsize\itshape, align=center}
]
  \node[box] (n1) {\textbf{Kap.~1--4}\\[5pt]Problem \&\\System\-landschaft};
  \node[box, right=0.5cm of n1] (n2) {\textbf{Kap.~5}\\[5pt]Anforde-\\rungsanalyse};
  \node[box, right=0.5cm of n2] (n3) {\textbf{Kap.~6--7}\\[5pt]Lösungsraum \&\\Bewertung};
  \node[box, right=0.5cm of n3] (n4) {\textbf{Kap.~8--9}\\[5pt]MVP \&\\Validierung};
  \node[box, right=0.5cm of n4] (n5) {\textbf{Kap.~10}\\[5pt]Fazit \&\\Ausblick};

  \draw[arr] (n1) -- (n2);
  \draw[arr] (n2) -- (n3);
  \draw[arr] (n3) -- (n4);
  \draw[arr] (n4) -- (n5);

  % DSR-Phasenbeschriftungen mit Klammerlinie
  % Relevance Cycle — Kap. 1–4
  \draw ([yshift=-6pt]n1.south west) -- ([yshift=-6pt]n1.south east);
  \node[phase] at ([yshift=-18pt]n1.south) {Relevance Cycle};

  % Rigor Cycle — Kap. 5–7 (überspannt n2 und n3)
  \draw ([yshift=-6pt]n2.south west) -- ([yshift=-6pt]n3.south east);
  \node[phase] at ($(n2.south)!0.5!(n3.south) + (0,-18pt)$) {Rigor Cycle};

  % Artefakt-Evaluation — Kap. 8–9
  \draw ([yshift=-6pt]n4.south west) -- ([yshift=-6pt]n4.south east);
  \node[phase] at ([yshift=-18pt]n4.south) {Artefakt-\\Evaluation};
\end{tikzpicture}
\vspace{0.2em}

\vspace{0.5em}\noindent\rule{\linewidth}{0.4pt}

% ── Strukturelle Entscheidungen ───────────────────────────────────────────────
{\large\bfseries Strukturelle Entscheidungen}\vspace{0.4em}

\renewcommand{\arraystretch}{1.2}
\begin{tabular}{@{}L{5.5cm} L{9.5cm}@{}}
\toprule
\textbf{Entscheidung} & \textbf{Begründung} \\
\midrule
DSR als Methodik & Zwei Artefakte im Zentrum: Scoring-Modell + MVP \\
Kap.\ 2.5 Entscheidungsmethoden im Software Engineering & Theoretische Basis für Kap.\ 7 --- Methodenwahl (AHP, Nutzwertanalyse etc.) wird hier fundiert \\
Kap.\ 4.4 Datenbereitstellung & Janine: explizit abzudeckendes Thema \\
Kap.\ 6.6/6.7 Referenzprojekte + Machbarkeitseinschätzung & Janines Kernkriterium: Umsetzbarkeit als Vorfilter vor dem Scoring \\
Kap.\ 7.1 Wahl der Bewertungsmethodik & Sachse: Warum dieses Bewertungsverfahren? --- explizit begründen \\
\bottomrule
\end{tabular}

\vspace{0.5em}\noindent\rule{\linewidth}{0.4pt}

% ── Zeitplan ─────────────────────────────────────────────────────────────────
{\large\bfseries Zeitplan}\vspace{0.4em}

\noindent Das Schreiben der Kapitel läuft \textbf{parallel zu jeder Phase} --- nicht erst am Ende.

\vspace{0.4em}
\begin{tabular}{@{}p{2.5cm} p{2.8cm} p{5.0cm} p{4.2cm}@{}}
\toprule
\textbf{Phase} & \textbf{Zeitraum} & \textbf{Schwerpunkt} & \textbf{Parallel schreiben} \\
\midrule
Orientierung      & 20. Apr -- 7. Mai  & Onboarding, Systemlandschaft, Stakeholder & Kap. 1, 3 + Literatur \\
Anforderungen     & 8. Mai -- 24. Mai  & Interviews, Requirements erheben          & Kap. 2, 4 \\
Analyse \& Modell & 25. Mai -- 14. Jun & Lösungsraum, Scoring-Modell               & Kap. 5, 6, 7 \\
MVP               & 15. Jun -- 30. Jun & MVP umsetzen und validieren               & Kap. 8, 9, 10 \\
Puffer            & 1. Jul -- 9. Jul   & Korrekturen, Feinschliff                  & --- \\
\midrule
\textbf{Abgabe}   & \textbf{10. Juli 2026} & & \\
\bottomrule
\end{tabular}

\end{document}
